\subsection{表示的特征标}

设 E 是 K 上有限维的左 A-模。E 的特征标或 E 的迹，记作 \(Tr_{E}\)，是线性型 \(a \mapsto \mathrm{Tr}(a_{\mathrm{E}})\)。设 \((e_{1}, \ldots, e_{n})\) 是 E 的一组基，\((e_{1}^{*}, \ldots, e_{n}^{*})\) 是它的对偶基。按定义，我们有关系式

\[
\mathrm{Tr} _ {\mathrm{E}} = \sum_ {i = 1} ^ {n} c _ {\mathrm{E}} (e _ {i}, e _ {i} ^ {*}).\tag{7}
\]

线性型 \(Tr_{E}\) 属于 \(\Theta_{\mathrm{E}}(\mathrm{A})\)。我们有 \(Tr_{E} = Tr_{E'}\)，这里 E 与 \(E'\) 是两个同构的 A-模。因此 \(Tr_{E}\) 只依赖于 E 的同构类。

设 \(\pi\) 是 A 在 K 上有限维向量空间 V 中的一个线性表示。\(\pi\) 的迹（有时也称 \(\pi\) 的特征标）定义为 A-模 \(V_{\pi}\) 的特征标，我们把它记作 \(\mathrm{Tr}_{\pi}\)。若 \((e_1, \ldots, e_n)\) 是 V 在 K 上的一组基，且 \((\pi_{ij}(a))\) 是 \(\pi(a)\) 关于这组基的矩阵，则我们有

\[
\operatorname{Tr} _ {\pi} (a) = \sum_ {i = 1} ^ {n} \pi_ {i i} (a).\tag{8}
\]

线性型 \(\mathrm{Tr}_{\pi}\) 属于 \(\Theta_{\pi}(\mathrm{A})\)。

命题 5. —— 设 S 是 K 上有限维的单 A-模，D 是 S 的交换子的反体，Z 是 D 的中心。下列性质等价：

\begin{enumerate}
\def\labelenumi{(\roman{enumi})}
\item
  S 的特征标 \(Tr_{S}\) 不为零。
\item
  存在元素 \(d \in \mathrm{D}\) 使得 \(\operatorname{Tr}_{\mathrm{D} / \mathrm{K}}(d) \neq 0\)。
\item
  \(\mathbf{Z}\) 是 \(\mathbf{K}\) 的可分扩张，且特征 \(p\)（\(\mathbf{K}\) 的特征）不整除次数 \([\mathrm{D}:\mathrm{Z}]\)。
\end{enumerate}

右 D-向量空间 S 是有限维的。设 \((e_{1},\ldots,e_{n})\) 是 S 在 D 上的一组基，u 是 \(\operatorname{End}_{\mathrm{D}}(\mathrm{S})\) 的一个元素。设 \((d_{ij})\) 是 u 关于基 \((e_{1},\ldots,e_{n})\) 的矩阵。我们有 \(u(e_{j})=\sum_{i=1}^{n}e_{i}d_{ij}\)（对 \(j\in[1,n]\)）。用 \(u_{K}\) 表示把 u 视为 K-向量空间 S 的自同态而得到的映射，用 \((u_{ij})\) 表示 \(u_{K}\) 关于 K-向量空间 S 的直和分解 \((e_{i}\mathrm{D})\) 的矩阵（II, §10, No.~5, 第 346 页）。我们有 \(\operatorname{Tr}(u_{\mathrm{K}})=\sum_{i}\operatorname{Tr}(u_{ii})\)。此外，\(u_{ii}\) 是 K-向量空间 \(e_{i}D\) 的自同态，它由 \(u_{ii}(e_{i}d)=e_{i}d_{ii}d\)（\(d\in D\)）定义，因此其迹等于 \(\operatorname{Tr}_{\mathrm{D}/\mathrm{K}}(d_{ii})\)。这样我们就证明了等式

\[
\mathrm{Tr} (u _ {\mathrm{K}}) = \mathrm{Tr} _ {\mathrm{D/K}} \bigl (\sum_ {i} d _ {i i} \bigr).\tag{9}
\]

由伯恩赛德定理（VIII, 第 83 页，命题 4 的推论 1），映射 \(a \mapsto a_{\mathrm{S}}\)（从 A 到 \(\mathrm{End}_{\mathrm{D}}(\mathrm{S})\) 的映射）是满射。性质 (i) 与 (ii) 的等价性于是由公式 (9) 得出。

若 K 的扩张 Z 是可分的，则存在元素 \(d \in Z\) 使得 \(\operatorname{Tr}_{\mathrm{Z/K}}(d) \neq 0\)（V, §8, No.~2, 第 49 页，命题 1 的推论）。此外，我们有 \(\operatorname{Tr}_{\mathrm{D/K}}(d) = [\mathrm{D} : \mathrm{Z}] \operatorname{Tr}_{\mathrm{Z/K}}(d)\)（III, §9, No.~4, 第 548 页，推论）；因此，若 p 不整除 {[}D : Z{]}，则我们有 \(\operatorname{Tr}_{\mathrm{D/K}}(d) \neq 0\)。这就证明了 (iii) 蕴含 (ii)。

若 \(\mathrm{Z}\) 不是 \(\mathrm{K}\) 的可分扩张，则 \(\operatorname{Tr}_{\mathrm{Z / K}}(x) = 0\)（对每个 \(x \in \mathrm{Z}\)），从而 \(\operatorname{Tr}_{\mathrm{D / K}}(d) = \operatorname{Tr}_{\mathrm{Z / K}}(\operatorname{Tr}_{\mathrm{D / Z}}(d)) = 0\)（对每个 \(d \in \mathrm{D}\)）。

现设 p 整除次数 \([D:Z]\)。于是它整除 D 在 Z 上的约化次数。对每个 \(d \in D\)，我们有 \(\mathrm{Tr}_{\mathrm{D}/\mathrm{Z}}(d) = 0\)（VIII, 第 344 页，注），从而 \(\mathrm{Tr}_{\mathrm{D}/\mathrm{K}}(d) = \mathrm{Tr}_{\mathrm{Z}/\mathrm{K}}(\mathrm{Tr}_{\mathrm{D}/\mathrm{Z}}(d)) = 0\)。这就完成了蕴含关系 (ii) \(\Rightarrow\) (iii) 的证明。

推论. —— 若体 \(\mathbf{K}\) 是完满的，则性质 (i)、(ii)、(iii) 均成立。

由于该体是完满的，K 的扩张 Z 是可分的（V, §15, No.~5, 第 125 页，定理 3）。性质 (iii) 于是由 VIII, 第 323 页的推论 3 得出。

命题 6. —— 设 \(\mathcal{S}_0\) 是 K 上有限维且迹非零的单 A-模的类所成的集合。线性型族 \((\mathrm{Tr}_{\mathrm{S}})_{\mathrm{S}\in \mathcal{S}_0}\) 在 K 上线性无关。

设 \(\mathrm{F}\) 是 \(\mathcal{S}_0\) 的一个有限子集，\(\mathrm{S}\) 是 \(\mathrm{F}\) 的一个元素。按假设，存在元素 \(a \in \mathrm{A}\) 使得 \(\operatorname{Tr}_{\mathrm{S}}(a) \neq 0\)。由命题 4 的推论 1（VIII, 第 83 页），存在元素 \(b \in \mathrm{A}\) 使得 \(b_{\mathrm{S}} = a_{\mathrm{S}}\)，且 \(b_{\mathrm{T}} = 0\)（对每个 \(\mathrm{T} \in \mathrm{F} - \{\mathrm{S}\}\)）。我们有 \(\operatorname{Tr}_{\mathrm{S}}(b) \neq 0\)，且 \(\operatorname{Tr}_{\mathrm{T}}(b) = 0\)（对 \(\mathrm{T} \in \mathrm{F} - \{\mathrm{S}\}\)）。于是族 \((\operatorname{Tr}_{\mathrm{S}})_{\mathrm{S} \in \mathrm{F}}\) 线性无关。命题 6 得证。

注. —— 设 \(S_{K}\) 是 K 上有限维单 A-模的类所成的集合。命题 6 也可由下述事实得到：诸 \(\Theta_{\mathrm{S}}(\mathrm{A})\)（S 取遍 \(S_{K}\)）之和是直和。

设

\[
0 \longrightarrow \mathrm{E} ^ {\prime} \longrightarrow \mathrm{E} \longrightarrow \mathrm{E} ^ {\prime \prime} \longrightarrow 0
\]

是 A-模的一个正合列，其中各模都在 K 上有限维。由 III, §9, No.~2, 第 543 页的命题 1，我们有 \(\mathrm{Tr}_{\mathrm{E}} = \mathrm{Tr}_{\mathrm{E}'} + \mathrm{Tr}_{\mathrm{E}''}\)。由格罗滕迪克群 \(\mathrm{R}_{\mathrm{K}}(\mathrm{A})\) 的定义（VIII, 第 194 页）及其泛性质（VIII, 第 186 页，命题 4），存在一个加群同态 \(\theta\)，从 \(\mathrm{R}_{\mathrm{K}}(\mathrm{A})\) 到 \(\mathrm{A}^*\)，它由关系式 \(\theta([\mathrm{E}]) = \mathrm{Tr}_{\mathrm{E}}\)（对每个 K 上有限维的 A-模 E 成立）刻画。特别地，E 的半单化的迹等于 E 的迹。由 \(\theta\) 我们导出一个 K-线性映射 \(\theta_{\mathrm{K}}: \mathrm{K} \otimes_{\mathbf{Z}} \mathrm{R}_{\mathrm{K}}(\mathrm{A}) \to \mathrm{A}^*\)。

推论. —— a) 设体 K 的特征为 0。同态 \(\theta\) 与 \(\theta_{K}\) 都是单射。两个 K 上有限维的半单 A-模同构，当且仅当它们的特征标相等。

\begin{enumerate}
\def\labelenumi{\alph{enumi})}
\setcounter{enumi}{1}
\tightlist
\item
  设体 K 是特征为非零素数 p 的完满体。那么同态 \(\theta_{K}\) 是单射，且 \(\theta\) 的核为 \(p\mathrm{R}_{\mathrm{K}}(\mathrm{A})\)。设 E 是 K 上有限维的半单 A-模。线性型 \(Tr_{E}\) 为零，当且仅当存在 A-模 F 使得 E 同构于 \(F^{p}\)。
\end{enumerate}

设 \(S_{K}\) 是 K 上有限维单 A-模的类所成的集合。诸元素 {[}S{]}（其中 S 取遍 \(S_{K}\)）构成 Z-模 \(\mathrm{R}_{\mathrm{K}}(\mathrm{A})\) 的一组基，因而诸元素 \(1 \otimes [S]\) 构成 K-向量空间 \(\mathrm{K} \otimes_{\mathbf{Z}} \mathrm{R}_{\mathrm{K}}(\mathrm{A})\) 的一组基（VIII, 第 195 页）。

设体 K 是完满的。由命题 6 以及 VIII, 第 383 页的推论，元素 \(\theta([S]) = \theta_{\mathrm{K}}(1 \otimes [S]) = \mathrm{Tr}_{\mathrm{S}}\)（\(A^{*}\) 中的元素）在 K 上线性无关。由此可知同态 \(\theta_{K}\) 是单射，且同态 \(\theta\) 的核由元素 \(\sum_{S \in S_{K}} n_{S}[S]\)（\(\mathrm{R}_{\mathrm{K}}(\mathrm{A})\) 的元素，满足 \(n_{S} \cdot 1_{K} = 0\) 对每个 \(S \in S_{K}\)）组成。因此 \(\theta\) 的核等于 \(p R_{\mathrm{K}}(\mathrm{A})\)，其中 p 是 K 的特征。特别地，若 K 的特征为 0，则同态 \(\theta\) 是单射。

a) 的最后一个断言由 VIII, 第 190 页的推论得出。

设 b) 的假设成立，并设 E 是 K 上有限维的半单 A-模。若存在 A-模 F 使 E 是 p 个同构于 F 的子模的直和，则模 F 是半单的且在 K 上有限维，并且我们有 \(Tr_{E} = p Tr_{F} = 0\)。反之，设 \(Tr_{E} = 0\)。我们有 \([E] \in p R_{K}(A)\)，故每个单模 \(S \in S_{K}\) 在 E 中的重数都是 p 的倍数（VIII, 第 190 页，命题 7）；我们把它写作 \(pn_{S}\)，其中 \(n_{S} \in N\)。族 ( \(n_{S}\) ) 具有有限支集。令 \(F = \oplus_{S \in S_{K}} S^{n_{S}}\)。于是我们有 {[}E{]} = {[}F\^{}\{p\}{]}，从而 E 与 F\^{}\{p\} 同构（VIII, 第 190 页，推论）。

设 \(\mathrm{E}\) 是 \(\mathrm{K}\) 上有限维的 A-模。设 \(a \in \mathrm{A}\)。用 \(\chi_{\mathrm{E}}(a; \mathrm{T})\) 表示 \(1_{\mathrm{E}} + T a_{\mathrm{E}}\)（\(\mathrm{K}[\mathrm{T}]\)-模 \(\mathrm{E}[\mathrm{T}] = \mathrm{K}[\mathrm{T}] \otimes_{\mathrm{K}} \mathrm{E}\) 的自同态）的行列式（III, §8, No.~11, 第 541 页）。我们有关系式

\[
\chi_ {\mathrm{E}} (a; \mathrm{T}) \equiv 1 + \mathrm{Tr} _ {\mathrm{E}} (a) \mathrm{T} \pmod {\mathrm{T} ^ {2} \mathrm{K} [ \mathrm{T} ]}\tag{10}
\]

（同前引文，公式 (49)）。由于这个多项式的常数项等于 1，它是一个可逆的形式幂级数（IV, §4, No.~4, 第 30 页）。此外，若

\[
0 \longrightarrow \mathrm{E} ^ {\prime} \longrightarrow \mathrm{E} \longrightarrow \mathrm{E} ^ {\prime \prime} \longrightarrow 0
\]

是 A-模的一个正合列，且各模都在 K 上有限维，则我们有 \(\chi_{\mathrm{E}}(a;\mathrm{T})=\chi_{\mathrm{E}^{\prime}}(a;\mathrm{T})\chi_{\mathrm{E}^{\prime\prime}}(a,\mathrm{T})\)（III, §8, No.~7, 第 534 页，公式 (31)）。由格罗滕迪克群 \(\mathrm{R}_{\mathrm{K}}(\mathrm{A})\) 的定义（VIII, 第 194 页）及其泛性质（VIII, 第 186 页，命题 4），存在唯一的同态 \(\chi_{a}\)，从群 \(\mathrm{R}_{\mathrm{K}}(\mathrm{A})\) 到乘法群 \(1+\mathrm{TK}[[\mathrm{T}]]\)，使得对每个 K 上有限维的 A-模 E 有 \(\chi_{a}([\mathrm{E}])=\chi_{\mathrm{E}}(a;\mathrm{T})\)。由公式 (10) 我们导出关系式

\[
\chi_ {a} (x) \equiv 1 + \theta (x) (a) \mathrm{T} \pmod {\mathrm{T} ^ {2} \mathrm{K} [ [ \mathrm{T} ] ]}\tag{11}
\]

其中 \(x\in \mathrm{R}_{\mathrm{K}}(\mathrm{A})\)，\(a\in \mathrm{A}\)。

若 E 与 \(E'\) 是两个 K 上有限维且半单化同构的 A-模，则我们有 \(\chi_{\mathrm{E}}(a;\mathrm{T})=\chi_{\mathrm{E}'}(a;\mathrm{T})\)。

定理 2. —— 设 \(\mathcal{A}\) 是 K-向量空间 A 的一个生成子集。同态 \(\chi_{\mathcal{A}}: \mathrm{R}_{\mathrm{K}}(\mathrm{A}) \to (1 + \mathrm{TK}[[\mathrm{T}]])^{\mathcal{A}}\) 由关系式 \(\chi_{\mathcal{A}}(x) = (\chi_a(x))_{a \in \mathcal{A}}\) 定义，它是单射。

设 \(x\) 是 \(\mathrm{R}_{\mathrm{K}}(\mathrm{A})\) 中满足 \(\chi_{\mathcal{A}}(x) = 1\) 的一个元素。由 (11)，我们有 \(\theta(x)(a) = 0\)（对每个 \(a \in \mathcal{A}\)），从而 \(\theta(x) = 0\)，因为 \(\theta(x)\) 是 A 上的 K-线性型而 \(\mathcal{A}\) 生成 K-向量空间 A。若 K 的特征为零，这就蕴含 \(x = 0\)（VIII, 第 384 页，命题 6 的推论），从而在这种情形下结论成立。以下设 K 的特征 \(p\) 非零。

我们先处理 K 是代数闭的情形。由前述及同前引文，此时我们有 \(x \in p \mathrm{R}_{\mathrm{K}}(\mathrm{A})\)。设 \(y \in \mathrm{R}_{\mathrm{K}}(\mathrm{A})\) 满足 x = py。对 A 的每个元素 a，我们有 \(\chi_{a}(y)^{p} = \chi_{a}(py) = \chi_{a}(x) = 1\)。于是 \((\chi_{a}(y) - 1)^{p} = 0\)，故 \(\chi_{a}(y) = 1\)，因为环 K{[}{[}T{]}{]} 是整环。所以 y 属于自同态 \(\chi_{A}\) 的核。由归纳法可知，x 属于 \(p^{n}\mathrm{R}_{\mathrm{K}}(\mathrm{A})\)（对每个整数 \(n \geqslant 1\)）。由于 \(\mathrm{R}_{\mathrm{K}}(\mathrm{A})\) 是自由 Z-模，这蕴含 x = 0，从而在这种情形下 \(\chi_{A}\) 是单射。

若不再假设 K 是代数闭的，我们取定 K 的一个代数闭包 \(\overline{K}\)，并考虑如下的群与群同态构成的图

\[
\begin{array}{c} \mathrm{R} _ {\mathrm{K}} (\mathrm{A}) \xrightarrow {\chi_ {\mathcal {A}}} (1 + \mathrm{TK} [ [ \mathrm{T} ] ]) ^ {\mathcal {A}} \\ \Biggl \downarrow^ {u} \qquad \qquad \qquad \qquad \Biggl \downarrow^ {i} \\ \mathrm{R} _ {\overline {{\mathrm{K}}}} (\mathrm{A} _ {(\overline {{\mathrm{K}}})}) \xrightarrow {\overline {{\chi}} _ {\mathcal {A}}} (1 + \mathrm{T} \overline {{\mathrm{K}}} [ [ \mathrm{T} ] ]) ^ {\mathcal {A}}, \end{array}\tag{12}
\]

其中 \(u\) 是由 K 到 \(\overline{\mathbf{K}}\) 的标量扩张导出的同态（VIII, 第 195 页），\(i\) 是典范单射，而 \(\overline{\chi}_{\mathcal{A}}\) 是同态 \(z \mapsto (\chi_{1 \otimes a}(z))_{a \in \mathcal{A}}\)。由 III, §9, No.~1, 第 542 页的公式 (12)，图 (12) 交换。由前述，同态 \(\overline{\chi}_{\mathcal{A}}\) 是单射。由于 \(u\) 是单射（VIII, 第 195 页，定理 1），同态 \(\chi_{\mathcal{A}}\) 是单射。

推论 1. —— 设 E 与 F 是 K 上有限维的半单 A-模，\(\mathcal{A}\) 是 K-向量空间 A 的一个生成子集。假设对每个 \(a \in \mathcal{A}\)，K-向量空间 E 与 F 的自同态 \(a_{E}\) 与 \(a_{F}\) 的特征多项式都相等。那么 A-模 E 与 F 同构。

设 \(a\) 是 \(\mathcal{A}\) 的一个元素。\(a_{\mathrm{E}}\) 与 \(a_{\mathrm{F}}\) 的特征多项式有相同的次数，故 E 的维数等于 F 的维数；我们把它记作 \(n\)。设 \(\mathrm{Pc}_{\mathrm{E}}(a;\mathrm{T})\) 是 \(a_{\mathrm{E}}\) 的特征多项式。在 \(\mathrm{K(T)}\) 中，我们有等式

\[
\chi_ {\mathrm{E}} (a; \mathrm{T}) = \det \left(1 + a _ {\mathrm{E}} \mathrm{T}\right) = (- \mathrm{T}) ^ {n} \det \left(\frac {- 1}{\mathrm{T}} - a _ {\mathrm{E}}\right) = (- \mathrm{T}) ^ {n} \operatorname{Pc} _ {\mathrm{E}} \left(a; \frac {- 1}{\mathrm{T}}\right),
\]

而 \(\chi_{\mathrm{F}}(a;\mathrm{T})\) 由一个类似的公式给出。由我们的假设，有 \(\chi_{\mathrm{E}}(a;\mathrm{T})=\chi_{\mathrm{F}}(a;\mathrm{T})\)。由定理 2，我们有 \([E]=[F]\)，这蕴含 E 与 F 同构（VIII, 第 190 页，命题 7 的推论）。

推论 2. —— 设 A 是 K 上有限次数的中心单代数。设 B 是半单 K-代数，f 与 g 是从 B 到 A 的代数同态，\(\mathcal{B}\) 是 K-向量空间 B 的一个生成子集。下列性质等价：

\begin{enumerate}
\def\labelenumi{(\roman{enumi})}
\item
  存在 A 的内自同构 \(\theta\) 使得 \(g = \theta \circ f\)。
\item
  对每个 \(b \in \mathcal{B}\)，有 \(\mathrm{Pc}_{\mathrm{A / K}}(f(b); \mathrm{X}) = \mathrm{Pc}_{\mathrm{A / K}}(g(b); \mathrm{X})\)。
\end{enumerate}

当 \(\mathrm{K}\) 的特征为零时，这些性质等价于下面的性质：

\begin{enumerate}
\def\labelenumi{(\roman{enumi})}
\setcounter{enumi}{2}
\tightlist
\item
  对每个 \(b \in \mathcal{B}\)，有 \(\mathrm{Tr}_{\mathrm{A / K}}(f(b)) = \mathrm{Tr}_{\mathrm{A / K}}(g(b))\)。
\end{enumerate}

把 A 的加群分别赋予外部合成律 \((b,a)\mapsto f(b)a\) 与 \((b,a)\mapsto g(b)a\)，将这样得到的左 B-模记作 M 与 N。由 VIII, 第 257 页的命题 1，性质 (i) 等价于 B-模 M 与 N 同构这一事实。按构造，与 B 的元素 b 相应的位似 \(b_{M}\) 是 A 中由 \(f(b)\) 作的左乘；因此我们有关系式

\[
\begin{array}{c} \mathrm {Pc_ {M}} (b; \mathrm{X}) = \mathrm {Pc_ {A / K}} (f (b); \mathrm{X}), \\ \mathrm{Tr} _ {\mathrm{M}} (b) = \mathrm{Tr} _ {\mathrm{A/K}} (f (b)) \end{array}
\]

以及把 M 换成 N、f 换成 g 后的两个类似关系式。我们知道（VIII, 第 386 页，推论 1），B-模 M 与 N 同构当且仅当 \(\mathrm{Pc}_{\mathrm{M}}(b;\mathrm{X})=\mathrm{Pc}_{\mathrm{N}}(b;\mathrm{X})\) 对每个 \(b\in B\) 成立。当体 K 的特征为 0 时，该关系式也等价于 \(\mathrm{Tr}_{\mathrm{M}}(b)=\mathrm{Tr}_{\mathrm{N}}(b)\) 对每个 \(b\in B\) 成立（VIII, 第 384 页，命题 6 的推论）。性质 (i) 与 (ii) 的等价性，以及当 K 的特征为 0 时 (i) 与 (iii) 的等价性，由此得出。

